\documentclass[10pt,a4paper]{amsart}
\usepackage{amsmath,amssymb,amsthm,graphicx}
% Retained minimal preamble for the Stacks Project chapter varieties.tex.
% Verbatim excerpt of the chapter preamble: Xy-pic setup (native preamble lines 14-18), the native theorem
% declarations with the mathematical macros (native preamble lines 154-217), and the project own no-op
% definitions of the reference/slogan/history display environments (native wrapper.tex lines 1-3) that the
% counted statement uses for its slogan. Copied mechanically; nothing added, removed or redefined.
% For commutative diagrams we use Xy-pic
\usepackage[all]{xy}

% We use 2cell for 2-commutative diagrams.
\xyoption{2cell}
% Theorem environments.
%
\theoremstyle{plain}
\newtheorem{theorem}[subsection]{Theorem}
\newtheorem{proposition}[subsection]{Proposition}
\newtheorem{lemma}[subsection]{Lemma}

\theoremstyle{definition}
\newtheorem{definition}[subsection]{Definition}
\newtheorem{example}[subsection]{Example}
\newtheorem{exercise}[subsection]{Exercise}
\newtheorem{situation}[subsection]{Situation}

\theoremstyle{remark}
\newtheorem{remark}[subsection]{Remark}
\newtheorem{remarks}[subsection]{Remarks}

\numberwithin{equation}{subsection}

% Macros
%
\def\lim{\mathop{\mathrm{lim}}\nolimits}
\def\colim{\mathop{\mathrm{colim}}\nolimits}
\def\Spec{\mathop{\mathrm{Spec}}}
\def\Hom{\mathop{\mathrm{Hom}}\nolimits}
\def\Ext{\mathop{\mathrm{Ext}}\nolimits}
\def\SheafHom{\mathop{\mathcal{H}\!\mathit{om}}\nolimits}
\def\SheafExt{\mathop{\mathcal{E}\!\mathit{xt}}\nolimits}
\def\Sch{\mathit{Sch}}
\def\Mor{\mathop{\mathrm{Mor}}\nolimits}
\def\Ob{\mathop{\mathrm{Ob}}\nolimits}
\def\Sh{\mathop{\mathit{Sh}}\nolimits}
\def\NL{\mathop{N\!L}\nolimits}
\def\CH{\mathop{\mathrm{CH}}\nolimits}
\def\proetale{{pro\text{-}\acute{e}tale}}
\def\etale{{\acute{e}tale}}
\def\QCoh{\mathit{QCoh}}
\def\Ker{\mathop{\mathrm{Ker}}}
\def\Im{\mathop{\mathrm{Im}}}
\def\Coker{\mathop{\mathrm{Coker}}}
\def\Coim{\mathop{\mathrm{Coim}}}

% Boxtimes
%
\DeclareMathSymbol{\boxtimes}{\mathbin}{AMSa}{"02}

%
% Macros for moduli stacks/spaces
%
\def\QCohstack{\mathcal{QC}\!\mathit{oh}}
\def\Cohstack{\mathcal{C}\!\mathit{oh}}
\def\Spacesstack{\mathcal{S}\!\mathit{paces}}
\def\Quotfunctor{\mathrm{Quot}}
\def\Hilbfunctor{\mathrm{Hilb}}
\def\Curvesstack{\mathcal{C}\!\mathit{urves}}
\def\Polarizedstack{\mathcal{P}\!\mathit{olarized}}
\def\Complexesstack{\mathcal{C}\!\mathit{omplexes}}
% \Pic is the operator that assigns to X its picard group, usage \Pic(X)
% \Picardstack_{X/B} denotes the Picard stack of X over B
% \Picardfunctor_{X/B} denotes the Picard functor of X over B
\def\Pic{\mathop{\mathrm{Pic}}\nolimits}
\def\Picardstack{\mathcal{P}\!\mathit{ic}}
\def\Picardfunctor{\mathrm{Pic}}
\def\Deformationcategory{\mathcal{D}\!\mathit{ef}}
\newenvironment{reference}{}{}
\newenvironment{slogan}{}{}
\newenvironment{history}{}{}
\usepackage{cleveref}
\crefname{originallemma}{original lemma}{original lemmas}
\Crefname{originallemma}{Original lemma}{Original lemmas}
\makeatletter
\newlabel{algebra-lemma-colimits-NL}{{07BQ}{source}{}{sourcepointer-1}{}}
\newlabel{algebra-lemma-colimits-NL@cref}{{[originallemma][0][]07BQ}{[1][0][]source}{}{}{}}
\makeatother
\crefname{originallemma}{original lemma}{original lemmas}
\Crefname{originallemma}{Original lemma}{Original lemmas}
\makeatletter
\newlabel{algebra-lemma-exact-sequence-NL}{{00S2}{source}{}{sourcepointer-2}{}}
\newlabel{algebra-lemma-exact-sequence-NL@cref}{{[originallemma][0][]00S2}{[1][0][]source}{}{}{}}
\makeatother
\crefname{originalsection}{original section}{original sections}
\Crefname{originalsection}{Original section}{Original sections}
\makeatletter
\newlabel{algebra-section-etale}{{00U0}{source}{}{sourcepointer-3}{}}
\newlabel{algebra-section-etale@cref}{{[originalsection][0][]00U0}{[1][0][]source}{}{}{}}
\makeatother
\crefname{originallemma}{original lemma}{original lemmas}
\Crefname{originallemma}{Original lemma}{Original lemmas}
\makeatletter
\newlabel{lemma-tangent-space-cotangent-space}{{0B2D}{source}{}{sourcepointer-4}{}}
\newlabel{lemma-tangent-space-cotangent-space@cref}{{[originallemma][0][]0B2D}{[1][0][]source}{}{}{}}
\makeatother
\crefname{originallemma}{original lemma}{original lemmas}
\Crefname{originallemma}{Original lemma}{Original lemmas}
\makeatletter
\newlabel{modules-lemma-stalk-module-differentials}{{08TE}{source}{}{sourcepointer-5}{}}
\newlabel{modules-lemma-stalk-module-differentials@cref}{{[originallemma][0][]08TE}{[1][0][]source}{}{}{}}
\makeatother
\crefname{originallemma}{original lemma}{original lemmas}
\Crefname{originallemma}{Original lemma}{Original lemmas}
\makeatletter
\newlabel{morphisms-lemma-base-change-differentials}{{01V0}{source}{}{sourcepointer-6}{}}
\newlabel{morphisms-lemma-base-change-differentials@cref}{{[originallemma][0][]01V0}{[1][0][]source}{}{}{}}
\makeatother
\title{The tangent space at a point with separable residue field extension as the dual of the cotangent space}
\author{The Stacks Project Authors}
\date{stacks-project:varieties.tex:local-snapshot-20261011; GNU FDL 1.2 or later}
\begin{document}
\maketitle
\noindent\textit{Editorial note.} All delivered mathematical text is the author's own, copied from the pinned Stacks Project chapter file varieties.tex; no statement or proof step was rewritten, reordered or removed.\par\medskip
\noindent\textit{Editorial note.} The statement ends with the author's own remark that the result holds more generally for a separable algebraic residue field extension, and the proof is carried out in that generality; both are reproduced as written.\par\medskip
\noindent\textit{Editorial note.} The counted spans carry 6 navigation stubs and no bibliography citation.\par\medskip
\noindent\textit{Editorial note.} The preamble is the chapter's retained minimal preamble (sources/stacks-varieties/retained-preamble.tex lines 1-77), a verbatim excerpt of the native preamble.\par\medskip
\noindent\textit{Editorial note.} The label lemma-tangent-space-rational-point is defined only in varieties.tex, so no cross-chapter restatement had to be weighed for this card.\par\medskip
\noindent\textit{Allowed standard prerequisites (editorial).}\begin{itemize}
\item Duality between tangent space and fibre of the cotangent module. Same chapter, Lemma lemma-tangent-space-cotangent-space, invoked at proof lines 3080-3081. Supplies the second isomorphism from the first.
\item The fibre of the relative differentials depends only on the fibre, and the stalk of the differentials on the local ring. Morphisms, Lemma morphisms-lemma-base-change-differentials and Modules, Lemma modules-lemma-stalk-module-differentials, invoked at proof lines 3082-3085. Supplies the reductions to the spectrum of the local ring of the fibre.
\item Transitivity exact sequence for the cotangent complex and its behaviour under filtered colimits. Algebra, Lemma algebra-lemma-exact-sequence-NL and Lemma algebra-lemma-colimits-NL, invoked at proof lines 3095-3101. Supplies the reduction to the vanishing statements for the residue field over k.
\item A finite separable field extension is etale. Algebra, Section algebra-section-etale, invoked at proof lines 3102-3104. Supplies the vanishing of the cotangent complex of the residue field over k.
\end{itemize}\medskip
\ifdefined\hypertarget\hypertarget{sourcepointer-1}{}\fi
\noindent\textit{Source reference.} Original lemma 07BQ: sources/stacks-chow/tags.txt, source lines 9350--9350. Pointer outside the counted span, invoked once for context; the load-bearing argument of the counted statement is the author's proof reproduced in full in the proof span.\par\medskip
\ifdefined\hypertarget\hypertarget{sourcepointer-2}{}\fi
\noindent\textit{Source reference.} Original lemma 00S2: sources/stacks-chow/tags.txt, source lines 1024--1024. Pointer outside the counted span, invoked once for context; the load-bearing argument of the counted statement is the author's proof reproduced in full in the proof span.\par\medskip
\ifdefined\hypertarget\hypertarget{sourcepointer-3}{}\fi
\noindent\textit{Source reference.} Original section 00U0: sources/stacks-chow/tags.txt, source lines 1094--1094. Pointer outside the counted span, invoked once for context; the load-bearing argument of the counted statement is the author's proof reproduced in full in the proof span.\par\medskip
\ifdefined\hypertarget\hypertarget{sourcepointer-4}{}\fi
\noindent\textit{Source reference.} Original lemma 0B2D: sources/stacks-varieties/pinned-original-source.tex, source lines 3014--3014. Pointer outside the counted span, invoked once for context; the load-bearing argument of the counted statement is the author's proof reproduced in full in the proof span.\par\medskip
\ifdefined\hypertarget\hypertarget{sourcepointer-5}{}\fi
\noindent\textit{Source reference.} Original lemma 08TE: sources/stacks-chow/tags.txt, source lines 11162--11162. Pointer outside the counted span, invoked once for context; the load-bearing argument of the counted statement is the author's proof reproduced in full in the proof span.\par\medskip
\ifdefined\hypertarget\hypertarget{sourcepointer-6}{}\fi
\noindent\textit{Source reference.} Original lemma 01V0: sources/stacks-chow/tags.txt, source lines 2429--2429. Pointer outside the counted span, invoked once for context; the load-bearing argument of the counted statement is the author's proof reproduced in full in the proof span.\par\medskip
\begin{lemma}
\label{lemma-tangent-space-rational-point}
Let $f : X \to S$ be a morphism of schemes.
Let $x \in X$ be a point and let $s = f(x) \in S$.
Assume that $\kappa(x) = \kappa(s)$. Then there are canonical isomorphisms
$$
\mathfrak m_x/(\mathfrak m_x^2 + \mathfrak m_s\mathcal{O}_{X, x})
=
\Omega_{X/S, x} \otimes_{\mathcal{O}_{X, x}} \kappa(x)
$$
and
$$
T_{X/S, x} =
\Hom_{\kappa(x)}(
\mathfrak m_x/(\mathfrak m_x^2 + \mathfrak m_s\mathcal{O}_{X, x}),
\kappa(x))
$$
This works more generally if $\kappa(x)/\kappa(s)$ is a separable
algebraic extension.
\end{lemma}

\begin{proof}
The second isomorphism follows from the first by
Lemma \ref{lemma-tangent-space-cotangent-space}.
For the first, we can replace $S$ by $s$ and $X$ by $X_s$, see
Morphisms, Lemma \ref{morphisms-lemma-base-change-differentials}.
We may also replace $X$ by the spectrum of $\mathcal{O}_{X, x}$, see
Modules, Lemma \ref{modules-lemma-stalk-module-differentials}.
Thus we have to show the following algebra fact: let
$(A, \mathfrak m, \kappa)$ be a local ring over a field $k$
such that $\kappa/k$ is separable algebraic. Then the canonical map
$$
\mathfrak m/\mathfrak m^2
\longrightarrow
\Omega_{A/k} \otimes \kappa
$$
is an isomorphism. Observe that
$\mathfrak m/\mathfrak m^2 = H_1(\NL_{\kappa/A})$. By
Algebra, Lemma \ref{algebra-lemma-exact-sequence-NL}
it suffices to show that $\Omega_{\kappa/k} = 0$ and
$H_1(\NL_{\kappa/k}) = 0$. Since $\kappa$ is the union of
its finite separable extensions in $k$ it suffices to prove
this when $\kappa$ is a finite separable extension of $k$
(Algebra, Lemma \ref{algebra-lemma-colimits-NL}).
In this case the ring map $k \to \kappa$ is \'etale
and hence $\NL_{\kappa/k} = 0$ (more or less by definition, see
Algebra, Section \ref{algebra-section-etale}).
\end{proof}

\end{document}
